\section{Sensitivity Analysis}

Revision~4 reported one design curve at one parameter set. Because $n^\star$ and $n_{\min}$ jointly set the capital ask, the recruitment target and the phase plan, the design requires a stated uncertainty on both. Revision~5 specifies the analysis rather than reporting a result, because the input ranges are open items (O-9 through O-13).

The required output is a ranked variance-contribution analysis \cite{saltelli2008} over the inputs of Theorem~\ref{thm:nodecount} ($F$, $h$, $z_{\beta^\star}$, $\sigma$, $\rho$, $\Lambda$, $c_d$, $k$, $A$) and Proposition~\ref{prop:radius} ($\tau$, $T_{\mathrm{route}}$, the in-node and speed quantile functions, $\eta$), with each input's range sourced and cited, and with $n^\star$ and $n_{\min}$ reported as intervals rather than point values. Three interactions are known in advance to be material and must be reported explicitly: $\rho$ against $\sigma$ (both act on the pooling term, in opposite directions); $\beta^\star$ against the basket-size distribution (Proposition~\ref{prop:orderfill} makes safety stock a function of basket composition, not only of demand variability); and in-node time against $\eta$ (both enter $n_{\min}$, one linearly through the radius and one directly, so a metro with poor packing efficiency and slow picking compounds).

The illustrative sweeps already computed give the order of magnitude at stake. Holding all else at the Revision~4 values, $n^\star$ moves from 9.1 to 13.6 as $\rho$ moves from 0 to 0.9; and under the illustrative time and speed distributions of Section~\ref{sec:structural}, $n_{\min}$ changes by a factor of 3.5 between Revision~4's radius formula and Proposition~\ref{prop:radius}. Neither range was visible in Revision~4. The pooling range moves the deployed node count directly; the radius range does not bind on the illustrative parameters but determines whether $n_{\min}$ or $n^\star$ is the binding constraint, which is not knowable until O-9 and O-10 close.
